\section{结论}

本文以统一图网络和资源—现金守恒为底座，按照信息集从全知、当日知到多人反馈依次建立MILP、MDP与动态博弈，完成三个问题和六关的建模、求解与验证。

\begin{enumerate}
  \item \textbf{确定天气：}第一、二关全局最优终值分别为10470元和12730元。最优性证书之外，逐日回放证明路线、补给、挖矿、现金和负重全部合法，并纠正第二关公开流水的一处水量错误。
  \item \textbf{仅知当天气象：}第三关稳健最优策略为54/54箱三日直达，在三档高温概率和各1024条支撑天气中零失败。第四关从567候选中选出240/240箱、B5挖矿阈值策略，基准总体均值10162.9元，失败率0.055\%。
  \item \textbf{多人交互：}第五关的非对称均衡终值为9535和9510元，双方单边偏离增益均为0。第六关证明机械复制单人挖矿策略会因共享边外部性几乎必败；推荐185/185箱的分路直达策略，基准总体均值7767.6元，失败率0.0167\%。
\end{enumerate}

最重要的管理含义是：天气不确定时应把“留足返程资源”写成显式阈值；多人环境下应优先消除同向边重合，不能被单人矿山高收益诱导。方法上的核心则是把模型输出、规则合法性、统计风险和均衡稳定性分开验证。由此，论文的六关结论既可复现，也清楚说明了各自能被信任到什么程度。

\begin{resultbox}{最终推荐}
第一关执行10470元最优计划；第二关执行12730元最优计划；第三关54/54箱直达；第四关240/240箱、B5阈值挖矿；第五关采用异路错峰纯策略均衡；第六关每人185/185箱并公开随机分配两条直达路线与一个错峰角色。
\end{resultbox}
